\scanpage{85}
% manuscript lines 1
$\partial$-\uline{Catégorie} $\mathscr H$

% manuscript lines 2–6
\deleted{catégorie additive \uline{graduée}, i.e. où les $\operatorname{Hom}(X,Y)$ sont des groupes \uline{gradués} (la composition des morphismes étant additive sur les degrés), satisfaisant aux axiomes qui \deleted{vont} suivent. On désigne par $\mathscr H_{(0)}$ la catégorie $\mathscr H$ où les hom sont réduits au degré $0$.}

% manuscript lines 7–8
\deleted{\uline{Axiome} $\uline{\partial\,1}$ Pour tout $X\in\operatorname{Ob}\mathscr H$, il existe un \deleted{$\unsure{T}(Y)=Y(1)$} \deleted{$\unsure{T(Y)|X(1)}$} $\in\operatorname{Ob}\mathscr H$ et un isomorphisme fonctoriel \inserted{en $X$}, (sur $\mathscr H_{(0)}$) \deleted{en $X$}}

% manuscript lines 9
\[
\deleted{\operatorname{Hom}_1(X,Y)\overset{\sim}{\longrightarrow}\operatorname{Hom}_0(X,Y(1))}
\]

% manuscript lines 10
\deleted{Donc il y a un objet canonique}

% manuscript lines 11
\[
\deleted{i_Y\in\operatorname{Hom}_1(\deleted{\unsure{T}}\unsure{X(1)|T(Y)|Y(1)},Y)}
\]

% manuscript lines 12–13
\deleted{ayant une propriété universelle. Évidt \unsure{$(T(Y),i_Y)$|$(Y(1),i_Y)$} est défini à $\mathscr H_{(0)}$-isomorphisme près.}

% manuscript lines 14
\deleted{objet $X(1)$ de $\mathscr H$ et un \uline{isomorphisme} \uline{de degré $1$}}

% manuscript lines 15
\[
\deleted{\alpha_X:X(1)\underset{1}{\longrightarrow}X}
\]

% manuscript lines 16
\deleted{et un objet $X(-1)$ de $\mathscr H$ et un \uline{isomorphisme de degré $-1$}}

% manuscript lines 17
\[
\deleted{\beta_X:X(-1)\underset{-1}{\longrightarrow}X}
\]

% manuscript lines 18
\deleted{On en conclut, pour tt $Y$, des isomorphismes}

% manuscript lines 19–20
\[
\mathrm I\quad\left\{\begin{aligned}
\deleted{\operatorname{Hom}_n(X,Y)}&\deleted{\overset{\sim}{\longrightarrow}\operatorname{Hom}_{n+1}(X(1),Y)} .
\quad [=\operatorname{Hom}_{n-1}(X,\unsure{Y(1)})]\\
\deleted{\operatorname{Hom}_n(Y,X(1))}&\deleted{\unsure{=|\simeq|\overset{\sim}{\longrightarrow}}\operatorname{Hom}_{n+1}(Y,X)}
\end{aligned}\right.
\]

% manuscript lines 21–22
\[
\mathrm {II}\quad\left\{\begin{aligned}
\operatorname{Hom}_n(X,Y)&\overset{\sim}{\longrightarrow}\operatorname{Hom}_{n-1}(X(-1),Y)\\
\operatorname{Hom}_n(Y,X(-1))&\overset{\sim}{\longrightarrow}\operatorname{Hom}_{n-1}(Y,X)\deleted{\unsure{1}}
\end{aligned}\right.
\]
% archive sheet number: 84

